% Appendix C: the bibliometric study in full, moved out of Section 3 in loop0009
% item 04 (review change 7). Sources: paper2/popularity.md (the method paragraph,
% the sensitivity check), numbers from paper2.section2_figures.numbers(), the
% timeline figure from paper2.section2_figures.fig_timeline, the table from
% make_tables.names_table. The citation-network figure (sec2_fig3) was replaced
% in the body by tables/communities.tex and is no longer in the paper.
\section{The bibliometric study}
\label{app:names}

This appendix gives the method behind Section~\ref{sec:names} in full, the
counts per name and when each name was in use.

\paragraph{Search.} We measured how much each name is used with OpenAlex
\citep{Priem2022}, on 29 September 2026. For each of the twelve problems we
searched titles and abstracts for its name and its common variants, and
restricted the search to the four fields in which the problems are studied:
Computer Science, Mathematics, Engineering and Decision Sciences. Without that
restriction several names are ordinary words: ``narrowness'' returns over a
million works, and ``edge separation'' returns aerodynamics.

\paragraph{Relevance.} Inside those fields a search hit is still not a work that
uses the name. OpenAlex's search stems words and ignores punctuation, so a
quoted phrase matches more than the phrase. ``Narrowness'' matches
``narrowing'' in term rewriting; ``node searching'' matches look-ups in
peer-to-peer networks; ``split bandwidth'' matches radar interferometry and 5G
spectrum assignment; ``path-width'' matches the widths of roads, tool paths and
datapaths. We therefore fetched all 2,271 hits with their titles, abstracts,
topics and venues, and kept a work only if it is about the problem in the sense
of the Linhares--Yanasse table. For the eleven smaller names the 700 hits were
judged one by one. For pathwidth, the 1,571 hits were filtered by a rule: keep a
work that uses the one-word spelling, that OpenAlex places in a theory subfield,
or that mentions graphs, treewidth or minors and not roads, vehicles,
trajectories or datapaths. We checked the rule by reading samples of what it
kept and what it dropped. Of the 2,271 hits, 1,593 are relevant: 1,213 for
pathwidth and 380 for the other eleven names. The two search-game names keep
27\% and 40\% of their hits.

\paragraph{Robustness.} The counts are lower bounds. A work can study a problem
without naming it in its abstract, OpenAlex has few abstracts before about
1990, and its citation counts run below Google Scholar's. The judgements were
made by one annotator with model assistance, record a yes or no per work and not
the reason, and have not been checked by a second annotator. As a second reading
we applied a mechanical rule: the literal phrase with a context term and no
exclusion term. It agrees with the judgements on 543 of the 700 judged works
(78\%). Under both readings pathwidth has more than three times as many works as
the other eleven names together (3.2 and 3.6 times), the three empty names stay
empty apart from one 5G false positive, and the search games lose most of their
hits. MOSP's rank over all years is the one conclusion that depends on the
reading: fourth under the judgements, fifth under the rule.

\begin{table}[t]
\centering\small
\caption{Citations do not track usage: the most-cited papers are cited far
more often than their problem's name is used. Works that use each problem's
name in the sense of the Linhares--Yanasse table, and citations of the paper
the table cites for it.
\emph{Relevant}: hits kept by the relevance check. \emph{Citations}: OpenAlex
\texttt{cited\_by\_count}, 29 September 2026. Kashiwabara and Fujisawa (1979)
has no OpenAlex record. $^{a}$OpenAlex returned one pathwidth record twice: the
counts are of records, as Figure~\ref{fig:names} draws them, and there are 1,212
distinct relevant pathwidth works.}
\label{tab:names}
\resizebox{\linewidth}{!}{% Generated by python -m paper2.latex.make_tables. Source: paper2.section2_figures.numbers() (OpenAlex caches, 2026-09-29).
% Do not edit by hand.
\begin{tabular}{@{}llrrp{4.6cm}r@{}}
\toprule
Problem & Discipline & Relevant & Hits & Reference in the 2002 table & Citations \\
\midrule
Pathwidth & graph theory & 1,213$^{a}$ & 1,571 & \citet{LY2002ref13} & 214 \\
Gate matrix layout & VLSI & 110 & 125 & \citet{LY2002ref6}; \citet{LY2002ref8} & 134 / 89 \\
PLA folding & VLSI & 68 & 74 & \citet{LY2002ref6} & 134 \\
MOSP & operations research & 58 & 60 & \citet{Yanasse1997b}; \citet{LY2002ref4} & 75 / 71 \\
Vertex separation & graph theory & 55 & 100 & \citet{LY2002ref13} & 214 \\
Edge search & graph theory & 38 & 95 & \citet{LY2002ref10} & 293 \\
Node search & graph theory & 34 & 127 & \citet{LY2002ref9} & 124 \\
One-dimensional logic & VLSI & 11 & 18 & \citet{LY2002ref7} & 108 \\
Interval thickness & graph theory & 6 & 10 & \citet{LY2002ref5} & -- \\
Narrowness & graph theory & 0 & 35 & \citet{LY2002ref11} & 44 \\
Split bandwidth & graph theory & 0 & 35 & \citet{LY2002ref12} & 19 \\
Edge separation & graph theory & 0 & 21 & \citet{LY2002ref14} & 77 \\
\midrule
All twelve & & 1,593 & 2,271 & & \\
\bottomrule
\end{tabular}
% check: other eleven = 380; pathwidth distinct = 1212; fetched 2026-09-29, citations 2026-09-29
}
\end{table}

\paragraph{Counts and citations.} Table~\ref{tab:names} gives the counts per
name. Citation counts do not track name usage. \citet{LY2002ref10} is the
most-cited paper in the table (293 citations), while only 38 works use ``edge
search'' in its sense: the paper is cited as a foundational result on graph
searching, not because people work on the edge search game under that name. The
split is starker for \citet{LY2002ref7}: 108 citations against 11 works using
``one-dimensional logic''. Split bandwidth is Fomin's own term in the paper the
table cites \citep{LY2002ref12}, and no one adopted it.

\begin{figure}[t]
\centering
\includegraphics[width=\linewidth]{sec2_fig2_timeline.pdf}
\caption{The VLSI names belong to the 1980s, the search games peak later and
fade, MOSP peaks in 2005--2014, and only pathwidth is still rising. The shade is
the name's rate, relevant works per million works in the same four fields in
each five-year period, as a fraction of that name's own peak rate; the numbers
are raw counts. Rows are ordered by the period of their peak. Counts cover
1970--2024, since 2025--2026 is not fully indexed, so row totals are below those
of Table~\ref{tab:names}. Before 1990 (hatched) OpenAlex has few abstracts.}
\label{fig:timeline}
\end{figure}

\paragraph{When each name was in use.} Figure~\ref{fig:timeline} shows the
counts over time. PLA folding and gate matrix layout peak in 1980--1989 and fall
to almost nothing after 2010, and one-dimensional logic has no relevant work
after 2009. MOSP first appears in the 1990s, peaks in 2005--2014 and is still
active. Pathwidth's rate has roughly doubled since the 1990s, from about 15 to
27 works per million, and its 311 relevant works in 2020--2024 are more than
the whole history of any other name. Rows with a few works per period move by
whole cells on one paper, and should be read as anecdote rather than trend.
